\chapter*{Bastion Events}\stepcounter{chapter}\phantomsection\addcontentsline{toc}{chapter}{Bastion Events}\label{chapter:BastionEvents}
\DndDropCapLine{T}{\entryfont he following Bastion Events provide the DM with ready-to-use scenarios for selected events from the Dungeon Master's Guide (2024 Edition). Each scenario is tailored to the Kingdom of Fife and incorporates characters, locations, or circumstances connected to the campaign's overarching storyline. The DM may use these scenarios as presented, adapt them to the party's circumstances, or use them as inspiration for additional Bastion Events.}
\section*{Rule Changes}\stepcounter{section}\phantomsection\addcontentsline{toc}{section}{Rule Changes}
{\entryfont The players' Bastion operates under circumstances that differ from those of typical frontier fortresses or wilderness strongholds. To reflect its location within Crail's bustling citadel, the following modifications apply to the standard Bastion rules in the Dungeon Master's Guide (2024 Edition), p. 334 ff.}
\subsection*{Double the Events}
{\entryfont During each Bastion Turn, two Bastion Events are rolled instead of one. This reflects the dynamic and politically charged nature of Crail's citadel, where court intrigues, merchant rivalries, and street-level rumours abound.}
\subsection*{Replacing Attack with Robbery}
{\entryfont Due to the bastion's secure location within Crail's citadel, large-scale attacks are not plausible. Instead, Robberies replace standard Attack Events. When a Robbery occurs:
\begin{itemize}
	\item Roll as normal for an Attack event to determine how many successes and failures the bastion has (based on Bastion Defender presence and traits).
	\item For each 1 rolled on the event dice, the bastion suffers a successful theft unless a Bastion Defender is present. Such a Defender is then removed from the bastion.
	\begin{itemize}
		\item If no Defender is present, the bastion's treasury is reduced by 1d6 × 100 GP per 1 rolled.
		\item If the die (1d6) also results in a 1, a valuable or magical item may be stolen:
		\begin{itemize}
			\item This could be a decorative item on display (e.g., enchanted artwork, ceremonial armour, fine jewellery).
			\item Or it could be a magic item stored in the bastion and not actively carried or worn by the players.
		\end{itemize}
	\end{itemize}
\end{itemize}
}
\begingroup
	\DndSetThemeColor[PhbTan]
	\begin{DndComment}{DM Annotation}
		The DM selects a suitable item or uses a random table, and may allow recovery attempts via urban investigation, seedy contacts, or market tracking.
	\end{DndComment}
\endgroup
\subsubsection*{Robberies Discourages Visitors}
{\entryfont If the bastion was subjected to a successful robbery, no friendly NPCs will visit the bastion during the next Bastion Turn. Word travels fast in Crail, and few wish to associate with a location rumoured to be unsecured or under scrutiny by the city's less savoury elements. During the next Bastion Turn, any randomly rolled \textbf{Extraordinary Opportunity}, \textbf{Friendly Visitors}, or \textbf{Guests} Event is replaced by an \textbf{All Is Well} Event.}

\section*{Special Events}\stepcounter{section}\phantomsection\addcontentsline{toc}{section}{Special Events}
\subsection*{Event Requirements}
{\entryfont Certain Bastion Events require specific conditions to be met before they can occur, such as completing a particular quest, encountering a character, or reaching a certain point in the campaign. These requirements, if applicable, are listed at the beginning of each event's description. Events whose requirements have not been met cannot occur.}

\clearpage

\section*{Mechanical Assistance}\stepcounter{section}\phantomsection\addcontentsline{toc}{section}{Mechanical Assistance}
{\vspace*{-0.5em}\ \ \ \footnotesize\texttt{\textsc{bastion event - guest}}}\vspace*{0.15em}\\
{\entryfont \textbf{Requirements:} The players must have visited Auchtertool and defeated the Dread-Witch Queen of Cellardyke.}\\\\
{\entryfont Following the defeat of the Dread-Witch Queen of Cellardyke, the Mecha-King of Auchtertool seeks to strengthen his kingdom's relationship with the adventurers who aided his people. As a gesture of gratitude, he dispatches a specialized construct to their Bastion, tasked with providing technical assistance and improving its facilities.

Unfortunately, while the construct possesses considerable mechanical expertise, its understanding of the needs and customs of living creatures is somewhat limited.}

\clearpage

\section*{Remnants of Grandeur}\stepcounter{section}\phantomsection\addcontentsline{toc}{section}{Remnants of Grandeur}
{\vspace*{-0.5em}\ \ \ \footnotesize\texttt{\textsc{bastion event - extraordinary opportunity}}}\vspace*{0.15em}\\
{\entryfont \hyperref[cha:BrandonGrennach]{\LinkFont{Brandon Grennach}}, a Loxodon noble and survivor of the fall of Dundee, approaches the party with a proposal. Wishing to reunite the surviving nobility of his fallen hometown, he hopes to organize an elegant gathering where old acquaintances can reconnect and briefly escape the hardships brought upon them by Zargothrax's invasion.

Grennach offers to arrange the invitations and festivities, provided the party agrees to host the gathering at their Bastion and contribute 500 GP toward the preparations. In return, he assures them that such an occasion would provide an excellent opportunity to establish connections with some of Fife's most influential noble families.}
\subsection*{Gathering Preparations}
{\entryfont Before the gathering begins, the players can improve the festivities by providing additional amenities, entertainment, or refreshments from their Bastion.

Each meaningful contribution grants 1 Preparation Point, up to a combined maximum of 4.}
\begin{DndTable}{lXr}
	Contribution		& Examples												& Points	\\
	Decorations			& Flowers from Garden, decorative banners				& +1		\\
	Refreshments		& Meals from a Kitchen or Soup Kitchen, fine beverages	& +1		\\
	Entertainment		& Performances, games, musical entertainment			& +1		\\
	Special Facilities	& Entertainment Salon, Circus Tent, Gaming Hall, etc.	& +1 per facility
\end{DndTable}
\subsection*{Optional: Auction}
{\entryfont During the gathering, the players may organize an auction to sell valuable possessions, rare curiosities, or unusual treasures from their own collection.

\paragraph*{Audience Interest} For each item offered, roll a d10 and apply a modifier based on the item's rarity, as shown in the Item Rarity table. For items without a defined rarity, the DM assigns an equivalent rarity based on their scarcity and desirability.}
\begin{DndTable}[header=Item Rarity]{Xc}
	Item Rarity							& Interest Adjustment	\\
	Common								& -2		\\
	Uncommon \& Rare					& 0			\\
	Very Rare, Legendary \& Artifact	& +2		\\
\end{DndTable}
\vspace*{-1.2\fontdimen6\font}
\begin{DndTable}[header=Audience Interest]{lXc}
	d10				& Reaction								& Modifier	\\
	3 or lower		& Little interest in the item			& -2		\\
	4 - 6			& A few curious nobles					& 0			\\
	7 - 9			& Several interested collectors			& +2		\\
	10 or higher	& Two wealthy nobles compete fiercely	& +5		\\
\end{DndTable}
\vfill\eject
{\noindent\entryfont\paragraph*{Final Selling Price} The DM determines each item's estimated value based on its rarity, condition, and desirability among the attending nobles.

A character acting as auctioneer makes a \textbf{Charisma (Persuasion) Check} for each item offered, applying the modifier determined by the audience's interest. Particularly impressive presentations or convincing stories about an item's history may grant Advantage on the Check. The result determines the final selling price as a percentage of the item's estimated value, as shown in the Auction Results table.}
\begin{DndTable}[header=Auction Results]{lXr}
	d20 Test		&& Final Selling Price		\\
	9 or lower		&& 50\% of expected value	\\
	10 - 13			&& 75\% of expected value	\\
	14 - 17			&& 100\% of expected value	\\
	18 - 21			&& 150\% of expected value	\\
	22 or higher	&& 200\% of expected value	\\
\end{DndTable}
\subsection*{Rewards}
{\entryfont Hosting the gathering grants the players one \textbf{Friendly Visitor Event}, one \textbf{Treasure Event}, and a lasting contact among the nobility in the form of Brandon Grennach. If the Treasure Event's d100 roll exceeds 74, the players may enlarge one of their Special Facilities for free instead of receiving the usual treasure.

Additional preparations attract more guests and increase the potential rewards, as shown in the Gathering Rewards table.}
\begin{DndTable}[header=Gathering Rewards]{Xrrr}
	Preparation	& Extra Visitors	& Magical Discovery	& Theft	\\
	0			& -					& -					& -		\\
	1			& 1					& -					& 5\%	\\
	2			& 1					& 25\%				& 10\%	\\
	3			& 2					& 50\%				& 15\%	\\
	4			& 3					& 75\%				& 20\%	\\
\end{DndTable}
{\entryfont For each additional visitor, resolve one Friendly Visitor Event. Roll separately for Magical Discovery and Theft, using the probabilities shown.

\paragraph*{Theft} As the gathering grows, opportunistic thieves may exploit the festivities. If a theft occurs, the players lose 2d6 × 100 GP or a trinket of equivalent value, as determined by the DM.}

\begin{tikzpicture}[remember picture, overlay]%
	\node[yshift=-1cm, anchor=south] at (current page.south) {\includegraphics[width=\paperwidth]{%
		images/backgrounds/Remnants_of_Grandheur%
	}};%
\end{tikzpicture}%